\PassOptionsToPackage{francais}{babel}

\documentclass[12pt,a4paper]{report}
\usepackage[utf8]{inputenc}
\usepackage[T1]{fontenc}
%\usepackage{caption}
\usepackage{tikz}
\usepackage{fancyhdr}
\usepackage{makeidx}
\usepackage{lmodern}
%\usepackage{wrapfig}
\usepackage{float}
%\usepackage[grey,utopia]{quotchap}   %pour le style d'affichage du chapitre
\usepackage[centertags]{amsmath}      %pour centrer les formules mathematiques
\usepackage{amsthm}                   %pour les formules mathematiques
\usepackage{cite}                     %pour les cite de la biblio
\def\baselinestretch{1}               %pour les inter ligne e 1
%\usepackage{ulem}                    %pour souligner le chapitre
%\usepackage[Conny]{fncychap}         %Sonny, Lenny, Glenn, Conny, Rejne, Bjarne
%\usepackage{hyperref}%rouge
%\usepackage[colorlinks]{hyperref}
\usepackage{fancyhdr}
%\usepackage{colortab}
\usepackage{epsfig}
\usepackage{anysize}
%\usepackage[Bjarne]{fncychap}%englais one
%\usepackage[Rejne]{fncychap}%choix(encdre le titre)
%\usepackage[Sonny]{fncychap}
%\usepackage[Conny]{fncychap}%choix2
%\usepackage[utopia]{quotchap}
%\usepackage[Glenn]{fncychap}%choix3
%\usepackage[Lenny]{fncychap}% encadrer le chapitre
\usepackage[grey,utopia]{quotchap}
\usepackage{graphicx}
%\usepackage[algonl,boxruled,lined, french]{algorithm2e}
\usepackage{niceframe}
\usepackage[a4paper]{geometry}
%\usepackage[a4paper,left=2cm,right=2cm,top=2cm,bottom=2cm]{geometry}
\geometry{hscale=0.80,vscale=0.80,centering}
\usepackage{rotating}
\usepackage{fancybox}
\usepackage[english,francais]{babel}
%\selectlanguage{francais}
%\usepackage{lipsum}
%\usepackage{multirow}
%\usepackage{soul}


%Pour l'index
\usepackage{makeidx}
\makeindex

%Pour la liste des abréviation
\usepackage[french]{nomencl}
\makenomenclature





%****************************************************************************
%****************************************************************************
%							DÉBUT DU DOCUMENT			
%****************************************************************************
%****************************************************************************

\begin{document}
	
%****************************************************************************
%					   	 Ajout de la page de garde
%****************************************************************************	


%============================================================================
%
%
%						PAGE DE GARDE
%
%
%
%============================================================================

\begin{titlepage}



%****************************************************************************
%					   	     L'entête de la page
%****************************************************************************	

\begin{center}
	\scshape{{ République Algérienne Démocratique et Populaire}}\\
	\scshape{{Ministère de l'Enseignement Supérieur et de la Recherche Scientifique}}\\
	\scshape{{\large Université Abderrahmane Mira de Béjaïa}}
	
	\begin{figure}[H]
		\centering
		\includegraphics[width=0.4\linewidth]{pagedegarde/Logo_Univ_Bejaia}
		\label{fig:Logo_Univ_Bejaia}
	\end{figure}
	\scshape{{Faculté des Sciences Exactes}}\\
	\scshape{{Département d'Informatique}}\\\medskip
	\scshape{{\large Mémoire de Master Professionnel}}\\\medskip
	\scshape{\large{Option: }{\Large Génie Logiciel}}
\end{center}
\vspace{1cm}


%****************************************************************************
%					   	     Le thème du mémoire
%****************************************************************************

\boxput*(0,1){\colorbox{white}{\textbf{{\large Thème}}}}
{\setlength{\fboxsep}{25pt}
	\fbox{\begin{minipage}{13.5cm}
			\begin{center}
				{\scshape{\Huge Intitulé de votre thème Intitulé de votre thème Intitulé de votre thème}}
			\end{center}
		\end{minipage}}
}

\vspace{1.2cm}


%****************************************************************************
%					   	 Les noms des membres du groupe
%****************************************************************************

\begin{center}
	
	\textit{Présenté par :}\\\bigskip
	
	
	\begin{tabular}{lcr}
		\scshape{ALICE Bob} & \scshape{ALICE Bob}
	\end{tabular}
	
\end{center}


%****************************************************************************
%					   	 Les noms des membres du jury
%****************************************************************************

\vspace{0.5cm}

\begin{center}
	\textit{Soutenu devant le jury composé de :}\\\bigskip 
\end{center}
\begin{tabular}{p{3cm}p{6.5cm}p{2cm}p{4cm}}
	\textit{Président}   & Dr \scshape{Alice} Bob & M.C.A &  \scshape{U. A/Mira Béjaïa}\\
	\textit{Rapporteur}  & Dr \scshape{Alice} Bob & M.C.A &  \scshape{U. A/Mira Béjaïa}\\
	\textit{Examinateur} & Dr \scshape{Alice} Bob & M.C.A &  \scshape{U. A/Mira Béjaïa}\\
	\textit{Examinatrice}& Mme \scshape{Alice} Bob & M.A.A & \scshape{U.  A/Mira Béjaïa}\\
\end{tabular}

\vspace{0.7cm}

\vfill

\begin{center}
	\cornersize{3}
	\ovalbox{\textbf{Promotion}  $ 2015-2016 $}
\end{center}

\end{titlepage}



%****************************************************************************
%						   Ajout des remerciment
%****************************************************************************

%\thispagestyle{empty}
%\chapter*{Remerciements}
%===================================================================

Nous tenons à exprimer notre sincère reconnaissance et notre profonde gratitude aux personnes suivantes qui ont aidé à la production de cet ouvrage :\\

À Monsieur \textbf{Charles LOOMIS}, Monsieur \textbf{Marc-Elian BEGIN} et Madame \textbf{Louise MERIFIELD}, les BigBoss de l'entreprise \textbf{SixSq}, pour la confiance qu'ils m'ont accordé afin de bénéficier de ce stage de fin d'études.\\

À Monsieur \textbf{Charles LOOMIS}, mon encadrant au sein de l'équipe de développement, pour l'intérêt qu'il m'a transmis au sujet de la solution Big Data au sein de la plateforme SlipStream et pour la compréhension globale relative à ce thème que je suis parvenue à développer tout au long de mon stage à \textbf{SixSq}.\\


À Madame \textbf{Lilia SFAXI}, mon tuteur de stage à l'Institut National des Sciences Appliquées et de Technologie, modèle de douce patience, pour son merveilleux soutien scientifique, moral et amical, pour son patient travail de reconsolidation des acquis cognitifs et pour ses généreux conseils concernant la présentation de cet ouvrage.\\


À tout le cadre enseignant de l'\textbf{INSAT} et surtout mes professeurs qui par leur engagement scientifique et éducatif, durant ces cinq années d'études, ont été pour moi une source d'inspiration.\\


Mes vifs remerciements s'adressent aux membres du jury, pour l'honneur qu'ils m'ont fait en examinant ce mémoire de fin d'études, soyez assurés de ma respectueuse considération.



\thispagestyle{empty}
\pagenumbering{roman}

%****************************************************************************
%				Configuration de la Table des Matières
%****************************************************************************

\renewcommand{\contentsname}{Table des Matières}
\addcontentsline{toc}{chapter}{Table des Matières} %
\rhead{\itshape\bf{\emph{Table des Matières}}} \pagestyle{fancy}
\renewcommand{\headrulewidth}{0.4pt}
\renewcommand{\footrulewidth}{0.4pt}
\lhead{\emph{}} \tableofcontents

\newpage


%****************************************************************************
%				Configuration de la Table des Figures
%****************************************************************************


\renewcommand{\contentsname}{Table des Figures}
\addcontentsline{toc}{chapter}{Table des Figures} %
\rhead{\itshape\bf{\emph{Table des Figures}}} \pagestyle{fancy}
\renewcommand{\headrulewidth}{0.4pt}
\renewcommand{\footrulewidth}{0.4pt}
\lhead{\emph{}}\listoffigures


\newpage


%****************************************************************************
%				Configuration de la Liste des Tableaux
%****************************************************************************


\renewcommand{\contentsname}{Liste des Tableaux}
\addcontentsline{toc}{chapter}{Liste des Tableaux} %
\rhead{\itshape\bf{\emph{Liste des Tableaux}}} \pagestyle{fancy}
\renewcommand{\headrulewidth}{0.4pt}
\renewcommand{\footrulewidth}{0.4pt}
\lhead{\emph{}}\listoftables

\newpage
%****************************************************************************
%				Configuration de la Liste des abbréviations
%****************************************************************************
%Options -> Configurer Texmaker -> Commandes : remplacer la commande Makeindex par
%makeindex %.nlo -s nomencl.ist -o %.nls -t %.nlg 
%Compiler 1fois avec pdflatex ou latex apres 1 fois makeindex et 1 fois avec pdflatex ou latex

\renewcommand{\contentsname}{}
\addcontentsline{toc}{chapter}{Liste des Abréviations} %
\rhead{\itshape\bf{\emph{}}} \pagestyle{fancy}
\renewcommand{\headrulewidth}{0.4pt}
\renewcommand{\footrulewidth}{0.4pt}
\lhead{\emph{}}\listoftables
\renewcommand{\nomname}{Liste des Abréviations}
\printnomenclature
\nomenclature{$GL$}{Génie Logiciel}
\nomenclature{$ASR$}{Administration et Sécurité des réseaux}
\nomenclature{$OM$}{Omar Mawloud}
\newpage


%****************************************************************************
%					   	 Configuration des chapitres
%****************************************************************************

\thispagestyle{empty}
\chead{}

\rhead{\itshape  \leftmark}
\lhead{}

\pagestyle{fancy}
\renewcommand{\headrulewidth}{0.4pt}
\renewcommand{\footrulewidth}{0.4pt}



%****************************************************************************
%			Ajout de l'introduction générale + les chapitres
%****************************************************************************
\thispagestyle{empty}
\pagenumbering{arabic}

\chapter*{\scshape{Introduction Générale}}
\addcontentsline{toc}{chapter}{Introduction Générale}





%****************************************************************************
%					   	     Début du chapitre
%****************************************************************************

\chapter{Présentation de l'organisme d'accueil}
\thispagestyle{empty}


\newpage
\section*{Introduction}


%****************************************************************************
%					   	     Début du chapitre
%****************************************************************************
\chapter{Étude de l'existant}
\thispagestyle{empty}

\newpage
\section*{Introduction}
\nomenclature{$ASR$}{Administration et Sécurité des réseaux}

\paragraph{}


%****************************************************************************
%					   	     Début du chapitre
%****************************************************************************
\chapter{\scshape{Conception et réalisation}]
\thispagestyle{empty}


\newpage
\section*{Introduction}


%****************************************************************************
%					   	     Début du chapitre
%****************************************************************************
\chapter{\scshape{Implémentation}}\thispagestyle{empty}


\newpage
\section*{Introduction}




%****************************************************************************
%					  Ajout de la conclusion
%****************************************************************************

\newpage
\chapter*{Conclusion Générale et Perspectives}





%****************************************************************************
%			               Configuration de l'Annexe
%****************************************************************************

\appendix

\renewcommand{\contentsname}{Annexes}
\addcontentsline{toc}{chapter}{Annexes}
\rhead{\itshape\bf{\emph{Annexes}}} \pagestyle{fancy}
\renewcommand{\headrulewidth}{0.4pt}
\renewcommand{\footrulewidth}{0.4pt}
\lhead{\emph{}}

\setcounter{figure}{0} 
\setcounter{table}{0}
\setcounter{footnote}{0}
\setcounter{equation}{0}
\pagestyle{fancy}
\fancyhf{}
\renewcommand{\chaptermark}[1]{\markboth{\MakeUppercase{#1 }}{}}
\renewcommand{\sectionmark}[1]{\markright{\thesection~ #1}}
\fancyhead[RO]{\bfseries\rightmark}
\fancyhead[LE]{\bfseries\leftmark}
\fancyfoot[RO]{\thepage}
\fancyfoot[LE]{\thepage}
\renewcommand{\headrulewidth}{0.5pt}
\renewcommand{\footrulewidth}{0pt}

\makeatletter
\renewcommand\thefigure{A.\arabic{figure}}
\renewcommand\thetable{A.\arabic{table}} 
\makeatother


%****************************************************************************
%				       Affichage de l'Annexe
%****************************************************************************


\chapter{Annexe}
\thispagestyle{empty}
\graphicspath{{Annexe1/figures/}}
%==========================================================================

%    Annexe

%===========================================================================

\newpage
\section{Annexe A}




%****************************************************************************
%				       Affichage de l'index
%****************************************************************************

\printindex



%****************************************************************************
%		  	     Affichage la liste des abréviation
%****************************************************************************

\printnomenclature




%****************************************************************************
%				     Ajout de la Bibliographie
%****************************************************************************

\bibliographystyle{ieeetr}
\bibliography{Biblio/bibliographie.bib}



%****************************************************************************
%				    	 Ajout du résumé
%****************************************************************************

\newpage
\thispagestyle{empty}

\setlength{\parindent}{24pt}
\setlength{\topmargin}{50mm}

\setlength{\headheight}{0cm} 
\setlength{\headsep}{0cm}
\setlength{\textheight}{28cm}
\setlength{\textwidth}{15cm} 
\setlength{\marginparsep}{0cm} 
\setlength{\marginparwidth}{0cm}
\setlength{\headheight}{0cm} 
\setlength{\footskip}{0cm}

\begin{minipage}[c]{\textwidth}
\textbf{Résumé}
\\\\
\footnotesize
\paragraph{} Lorem ipsum \ldots \\
\underline{Mots clés :} xx, xxx, xx.
\end{minipage}
\\[1.5cm]

\textbf{Abstract}
\\\\

\begin{minipage}[c]{\textwidth}
\footnotesize

\paragraph{}Lorem ipsum \ldots \\

\underline{Keywords :} xx, xxx, xx.
\end{minipage}





\end{document}
